\documentclass[11pt,a4paper]{article}

\usepackage[utf8]{inputenc}
\usepackage[T1]{fontenc}
\usepackage{mathptmx}
\usepackage[margin=0.7in]{geometry}
\usepackage{enumitem}
\usepackage{hyperref}
\usepackage{titlesec}
\usepackage{xcolor}

\pagestyle{empty}
\setlength{\parindent}{0pt}
\setlist[itemize]{leftmargin=*, noitemsep, topsep=2pt}

\titleformat{\section}{\Large\bfseries}{}{0em}{}[\titlerule]
\titlespacing*{\section}{0pt}{10pt}{4pt}

\titleformat{\subsection}[runin]{\bfseries}{}{0em}{}
\titlespacing*{\subsection}{0pt}{6pt}{0pt}

\definecolor{linkcolor}{RGB}{0,102,204}
\hypersetup{colorlinks=true, urlcolor=linkcolor}

\newcommand{\highlight}[1]{\textbf{#1}}

\begin{document}

\begin{center}
    {\Huge \textbf{ERIC MAYEUX}}\\
    \vspace{3pt}
    {\large Directeur Qualit\'{e} \& Technologies | \highlight{Leadership Agile} | \highlight{Gestion d'\'{E}quipes TI}}\\
    \vspace{3pt}
    Montr\'{e}al, QC, Canada\\
    \href{mailto:mayeux.e@gmail.com}{mayeux.e@gmail.com} | 438-884-7645
\end{center}

\vspace{6pt}

\section*{Profil Professionnel}
\noindent
Directeur Qualit\'{e} et Technologies avec plus de 10 ans d'exp\'{e}rience en TI, sp\'{e}cialis\'{e} en \highlight{gestion d'\'{e}quipes multidisciplinaires}, \highlight{d\'{e}veloppement logiciel} et transformation des \highlight{syst\`{e}mes}. Reconnu pour son \highlight{leadership influent}, son \highlight{sens strat\'{e}gique} et sa \highlight{d\'{e}termination} \`{a} livrer des r\'{e}sultats. Leader mobilisant par la \highlight{collaboration intersectorielle}, la \highlight{discipline} et l'\highlight{esprit d'\'{e}quipe} dans un contexte \highlight{Agile} (\highlight{Scrum}, \highlight{SAFe}, \highlight{Kanban}). Ma\^{i}trise du \highlight{soutien op\'{e}rationnel}, des \highlight{technologies} d'automatisation et de l'\highlight{innovation} appliqu\'{e}e. Exp\'{e}rience internationale (Canada, France, Tha\"{i}lande) en milieu bancaire.

\section*{Comp\'{e}tences Cl\'{e}s}

\textbf{Leadership \& Gestion:} \highlight{Gestion d'\'{e}quipes multidisciplinaires} (5+ ressources), \highlight{Gestion des ressources humaines}, Strat\'{e}gie et gouvernance, KPIs qualit\'{e}, Mentorat, Recrutement, Formation, \highlight{Accompagnement au changement}

\textbf{M\'{e}thodologies Agile:} \highlight{Scrum}, \highlight{SAFe}, \highlight{Kanban}, DevOps, PI Planning, TDD, BDD

\textbf{D\'{e}veloppement Logiciel \& Technologies:} Java, Python, JavaScript, Shell, SQL, PL/SQL, REST, SOAP, SOA, Microservices, Docker, Kubernetes, Jenkins, GitHub

\textbf{Automatisation \& Innovation:} Selenium, Playwright, RestAssured, JUnit, Gatling, K6, JMeter, IA appliqu\'{e}e (Crew AI, Playwright MCP), RPA

\textbf{Syst\`{e}mes \& Monitoring:} Splunk, Datadog, Jira, Confluence, Oracle, Sybase, \highlight{SLO/SLA}

\section*{R\'{e}alisations Cl\'{e}s}
\begin{itemize}
    \item Diriger et encadrer une \highlight{\'{e}quipe multidisciplinaire} de 5+ professionnels en contexte \highlight{agile SAFe}
    \item Piloter la mise en place de la Pratique Qualit\'{e} et de gouvernances p\'{e}rennes
    \item Conduire la transformation technologique : POC innovants, int\'{e}gration IA, organisation de hackathon
    \item Renforcer la \highlight{collaboration intersectorielle} entre squads, POs et parties prenantes
    \item Concevoir et d\'{e}ployer des frameworks d'automatisation et pipelines CI/CD/CT
    \item Contribuer au \highlight{d\'{e}veloppement logiciel} : applications m\'{e}tier, webservices, bases de donn\'{e}es
\end{itemize}

\section*{Exp\'{e}rience Professionnelle}

\subsection*{Int\'{e}grateur S\'{e}nior Qualit\'{e} | Senior Quality Integrator} \hfill \textit{F\'{e}v. 2025 -- Pr\'{e}sent}\\
\textbf{Banque Nationale du Canada (BNC)} -- Montr\'{e}al, QC
\begin{itemize}
    \item Piloter la \highlight{strat\'{e}gie} de test et la Pratique Qualit\'{e} pour un train SAFe multi-\'{e}quipes
    \item Encadrer, recruter et mentorer une \highlight{\'{e}quipe multidisciplinaire} de 5+ QA et int\'{e}grateurs
    \item Assurer la \highlight{collaboration intersectorielle} avec les squads locales et internationales (Tha\"{i}lande)
    \item D\'{e}finir des \highlight{SLO/SLA} et assurer la \highlight{r\'{e}solution rapide des incidents} en environnement critique
    \item Cr\'{e}er des tableaux de bord et m\'{e}triques de qualit\'{e} (Datadog, Splunk, Jira)
    \item Conduire des initiatives d'\highlight{innovation} : int\'{e}gration d'outils IA, hackathon de test
    \item Mettre en place des pipelines CI/CD/CT pour l'ex\'{e}cution continue des tests
\end{itemize}

\subsection*{Int\'{e}grateur Qualit\'{e} -- SDET | Quality Integrator} \hfill \textit{Jan. 2022 -- F\'{e}v. 2025}\\
\textbf{Banque Nationale du Canada (BNC)} via Groupe Astek -- Montr\'{e}al, QC
\begin{itemize}
    \item Concevoir et ex\'{e}cuter la \highlight{strat\'{e}gie} de test en environnement \highlight{Agile} (Gestion de Patrimoine \& ECM)
    \item Mettre en place des m\'{e}triques de qualit\'{e} : pyramide de test, couverture, tests de charge
    \item D\'{e}velopper et maintenir des \highlight{syst\`{e}mes} d'automatisation (RestAssured, Selenium, Gatling, K6)
\end{itemize}

\subsection*{QA Automaticien | QA Automation Engineer} \hfill \textit{Juin 2021 -- D\'{e}c. 2021}\\
\textbf{AODocs} -- Paris, France
\begin{itemize}
    \item Automatiser les tests de non-r\'{e}gression (Selenium, Java, JUnit 5, Karate)
\end{itemize}

\subsection*{Automaticien CI/CD | DevOps Automation Engineer} \hfill \textit{Jan. 2020 -- Juin 2021}\\
\textbf{ENGIE DIGITAL} via Groupe H\'{e}nix -- Paris, France
\begin{itemize}
    \item Piloter la cr\'{e}ation et la \highlight{maintenance} de la cha\^{i}ne CI/CD (Jenkins, Docker)
    \item Administrer les \highlight{syst\`{e}mes} et outils DevOps pour plusieurs projets TI
    \item Assurer la \highlight{gestion} op\'{e}rationnelle et le pilotage de projets via Jira (KPIs)
\end{itemize}

\subsection*{Product Owner} \hfill \textit{Juin 2019 -- Sept. 2019}\\
\textbf{H\'{e}nix} -- Montrouge, France
\begin{itemize}
    \item Piloter la recette et l'int\'{e}gration de plugins pour le Centre de Service Thales (8~500 utilisateurs Jira)
    \item G\'{e}rer les relations avec les \highlight{fournisseurs externes} et le support client SaaS
\end{itemize}

\subsection*{D\'{e}veloppeur Logiciel | Software Developer} \hfill \textit{Mai 2017 -- Mai 2019}\\
\textbf{MEDIAPOST} via Groupe H\'{e}nix -- Paris, France
\begin{itemize}
    \item Concevoir et \highlight{d\'{e}velopper} des webservices SOAP et REST en architecture SOA (BPEL)
    \item \highlight{D\'{e}velopper} des scripts de bases de donn\'{e}es (proc\'{e}dures stock\'{e}es PL/SQL, Sybase)
    \item Cr\'{e}er une application interne compl\`{e}te de \highlight{gestion} de donn\'{e}es comptables
\end{itemize}

\subsection*{Consultant QA \& Automatisation} \hfill \textit{Avril 2016 -- Avril 2017}\\
\textbf{ENGIE IT} via Groupe H\'{e}nix -- Saint-Ouen, France
\begin{itemize}
    \item Assurer la \highlight{gestion} op\'{e}rationnelle et la supervision des environnements de production
    \item Piloter les ressources sur plusieurs plateformes et la relation avec les \highlight{fournisseurs externes}
\end{itemize}

\section*{Certifications}

\textbf{PMP -- Project Management Professional} (en cours) \hfill \textit{2025}\\
\textbf{Certified Scrum Master (CSM)} -- Scrum Alliance \hfill \textit{2024}\\
\textbf{Jira ACP-600} -- Project Administration in Jira Server \hfill \textit{2019}\\
\textbf{Cursus Automaticien et DevOps} -- AFCEPF \hfill \textit{2019}\\
Selenium, SoapUI, NeoLoad, Docker, Kubernetes, Jenkins, Ansible\\
\textbf{Architecture logiciel \& Analyste informaticien -- Certifi\'{e} Master II} -- AFCEPF \hfill \textit{2017}

\section*{Formation}

\textbf{Licence en Gestion} -- \'{E}cole Sup\'{e}rieure de Gestion (E.S.G), Paris \hfill \textit{2011}\\
\textbf{D.P.E.C.F} -- Dipl\^{o}me Pr\'{e}paratoire \`{a} l'Expertise Comptable -- Advancia, Paris \hfill \textit{2007}

\section*{Langues}

\textbf{Fran\c{c}ais:} Langue maternelle \quad | \quad \textbf{Anglais:} Avanc\'{e} (Professional Working Proficiency)

\end{document}
